\chapter{Limitations and Future Work}
\label{ch:limitations}

Each successful stage of the project narrowed the remaining uncertainty. The simulator can now
configure, train, replay and evaluate substantially different control tasks, and the evaluated
policies have measurable operating ranges. The results do not establish a flight-ready
controller, superiority of PPO or robustness to every option visible in the UI. The remaining
work follows from three boundaries: the world being simulated, the experiments performed inside
it, and the software used to preserve them.

\section{Boundary of the physical model}

The physical model became more complete as new components required it, but it remains a
six-degree-of-freedom rigid body with configurable low-order forces. It does not include
propellant slosh, elastic structure, detailed engine transients, plume--ground interaction,
compressible or transonic aerodynamics, resolved grid-fin lattice flow, sensor noise or
realistic turbulent wind. Pressure-dependent thrust and specific impulse are approximations,
and landing contact is resolved by Unity colliders and a general-purpose solver. No flight data
or independent high-fidelity simulator validates the trajectory as a whole.

The measured policy performance is consequently valid inside the documented simulator. The
next useful physics step is not to add every missing effect at once. It is to determine which
current approximation changes the control outcome. Analytical force cases, published
aerodynamic data and comparison with a higher-fidelity tool could quantify coefficient
sensitivity before additional complexity is introduced.

\section{Boundary of the experiments}

Every completed result uses one principal training seed, one network architecture and one PPO
configuration. Wilson intervals describe repeated episodes for a fixed model; they do not
measure how much performance would vary across independently trained policies. Using one common
evaluation seed makes the five-band studies repeatable, but still samples only one deterministic
set of initial conditions.

The two curriculum tasks also fail differently at their hardest boundary. Tracking falls from
84\% success at $d=0.75$ to zero at $d=1$ after training stopped short of a stable full-difficulty
distribution. L11 landing trained at $d=1$ but achieved only 36\% there during deterministic
evaluation. These results lead to different next experiments. Tracking should be evaluated more
densely between $d=0.75$ and $d=1$ and trained with explicit probability mass at the boundary.
Landing should keep the final checkpoint and seeds fixed while changing only deterministic
versus stochastic inference.

If that landing comparison reproduces the training--evaluation gap, the next design change
should target the action interface rather than the reward. Engine enable is physically binary
but is currently produced by a thresholded continuous output. A discrete branch, or a deliberate
hybrid action space, would remove the dead band between the on and off thresholds. Only after
that isolated test should the main runs be repeated across several trainer and evaluator seeds.

Fixed hover now has a complete 250-episode evaluation, but still lacks a predeclared binary
criterion. The new policy corrects the earlier vertical drift and consistently damps velocity
and attitude, yet converges to a mean final planar offset of \SI{4.77}{m}. The next hover
experiment should define its acceptance radius and hold time before execution, then determine
whether the offset comes from reward balance, observations or the learned equilibrium.

The reward system has a related boundary. Configurability makes a specification inspectable but
does not make it correct. A future validator could estimate the maximum cumulative contribution
of every continuous term and warn when waiting can dominate a terminal outcome. Scripted
adversarial policies---prolonged hovering, engine-off descent and deliberate early
termination---could test these economics before an expensive training run begins.

\section{Boundary of the software product}

Architecture had to evolve while scenes, prefabs, sessions and ONNX policies already existed.
The current schema checks protect known interfaces, but they cannot guarantee migration of every
legacy artifact. Some agent partial classes remain large, not every UI path has an automated
test, and the final public commit is not yet frozen.

A command-line experiment runner would be the most useful complement to the interactive UI. It
could queue matrices of seeds and configurations without manual setup and make larger comparisons
less error-prone. Automated edit- and play-mode tests, a dependency lock, a small known-answer
evaluation and an explicit artifact migration tool would then make the public release easier to
verify and maintain.

\section{A deliberate order for future experiments}

The simulator exposes many possible extensions, but repeating the original scope problem would
make their results difficult to interpret. The next work should therefore proceed in the order
in which each experiment unlocks the next:

\begin{enumerate}
  \item define fixed-hover acceptance limits, diagnose its planar offset and isolate the L11
  deterministic-inference gap;
  \item repeat the principal policies across several seeds once their protocols are stable;
  \item compare selected vehicle configurations on one frozen task, retraining whenever the
  policy interface changes;
  \item measure a no-wind policy under controlled wind, then compare it with a policy trained
  through domain randomization;
  \item inject one actuator fault at a time with fixed onset, duration and severity before
  attempting combined faults; and
  \item treat tower-arm capture as a separate task with its own contact model and evaluation.
\end{enumerate}

This order keeps the core claim modest: the current thesis demonstrates the configurable
workflow and three control tasks. Wind, faults, vehicle comparisons and tower capture remain the
next chapters of the product rather than unevaluated additions to the present result.

\section{Final thesis completion checklist}

\begin{thesistodo}
  Run the deterministic/stochastic L11 inference diagnostic; declare hover acceptance limits
  and diagnose its planar offset;
  freeze and publish the repository revision; add UI/task screenshots and seed-labelled videos;
  translate the approved text to Slovene; insert the official STUDIS description, mentor titles,
  acknowledgements and student email; update both abstracts after the final accepted results;
  and have the mentor approve the final title and scope.
\end{thesistodo}
